\documentclass[11pt,a4paper]{article}
\usepackage[margin=24mm]{geometry}
\usepackage{amsmath,amssymb,amsthm,booktabs,array}
\input{glyphtounicode}
\pdfgentounicode=1
\usepackage[hidelinks]{hyperref}
\hypersetup{hypertexnames=false}
\newcommand{\pos}[1]{[#1]_+}
\newcommand{\dd}{\,\mathrm d}
\newcommand{\kc}{k_c}
\newcommand{\Bud}{\mathcal B}
\newcommand{\He}{\mathcal H}
\newtheorem{proposition}{Proposition}
\newtheorem{theorem}{Theorem}
\title{AS3: asymmetric initial lanes, useful finite advice changes\newline
and the complete gradient reference}
\author{Elementary VSL lane-drop task review}
\date{19 October 2026}
\begin{document}
\maketitle
\noindent\textbf{Status.} AS1 and AS2 are user-accepted. This completes AS3
at review/reference scope and returns to the original Step 3 user closure
gate. No traffic rollout, numerical derivative, AD or optimization is run.
No new traffic law or numerical smoothing is selected.

\paragraph{Result and limitation.}
The accepted hidden A/B split can matter for complete-TST advice, not merely
for local receiving: a conditional existence range of the selected model
has opposite utility for the same finite slowdown. A nonzero common-fastest
range also exists. The positive range may require large forced work, weak
mismatch and fast recovery; it is not an admitted physical parameter range.
The complete regular-event sensitivity below is a reference for later
verification, not proof of numerical gradient accuracy or usefulness.

\section{Accepted traffic and information remain fixed}
Use the
\href{vsl_as1_asymmetric_forward_contract.pdf}{accepted AS1 forward contract}
and \href{vsl_as2_information_and_history.pdf}{accepted AS2 observation contract}.
E ends, A is its adjacent continuing lane and B the far lane. The initial
uniform footprint is $[0,\ell]$, $L=L_p+b$ and
\begin{equation}
 e_0,a_0,b_0\geq0,\quad r=e_0+a_0,\quad
 \rho=e_0+a_0+b_0,\quad M=\ell\rho.
 \label{eq:initial}
\end{equation}
Each true initial state is fixed as advice $v\in[v_-,v_+]$ varies.
So are geometry, physical coefficients, the W-F anchor and $m=1$ veh-eq.
Retain independent lane loads, neutral initial terminal/recovery histories,
no postdecision arrivals, finite storage and choke completion. In particular,
\begin{equation}
 v_+r<R\leq C(v_-),\quad v_+b_0\leq R_B\leq C(v_-),\quad
 M<k_jb,\quad c_0-\beta(1+v_+/w)>0,
 \label{eq:domain}
\end{equation}
with all other AS1 bounds intact. Write $K=wk_j$,
$C(v)=Kv/(w+v)$, $\kc(v)=K/(w+v)$, $U(g)=wg/(K-g)$ and
$c_0=1+\delta\geq1$, $\chi=\beta/[C(v_+)v_+]$ fixed once.

Level 2 receives E and continuing totals and every permitted history/inference;
Level 3 additionally knows A/B separately. The compact peer knows one extra
initial adjacent mass and exactly matches Level 3. No realized $r$, future
E/service trace or pair label is leaked into Level 2. AS2's stronger matched
aggregate/E history is retained, and identifying histories remain admissible.

\section{A genuinely asymmetric analytical reference family}
\subsection{The same E and continuing total, with unequal adjacent loading}
For arbitrary positive scales $u,\kappa,h$, choose
\begin{equation}
 \begin{gathered}
 v_-=u,\ v_+=2u,\ w=2u,\ k_j=3\kappa,\ \gamma=2,\quad
 R=39\kappa u/20,\ R_B=\kappa u,\\
 \nu=\ell/h=\log(6/5),\quad L_p=2\ell,\quad b=\ell,\quad L=3\ell,\\
 \mathcal L:(e_0,a_0,b_0)=(3\kappa/5,\kappa/5,7\kappa/20),\\
 \mathcal H:(e_0,a_0,b_0)=(3\kappa/5,7\kappa/20,\kappa/5).
 \end{gathered}\label{eq:family}
\end{equation}
Here $\rho=23\kappa/20$ is common, $r_{\mathcal L}=4\kappa/5$,
$r_{\mathcal H}=19\kappa/20$, and $M/(k_jb)=23/60<1$.
The worst ordinary load margin is $R-v_+r_{\mathcal H}=\kappa u/20$.
AS2's same lower-speed loading construction reaches both initial states
with identical full permitted histories, including E history.

The family fixes initial state/geometry; it does not fix one work cost for
all deductions. AS2 illustrated $c_0=11/10,\beta=1/100$. The action existence
proof below chooses another \emph{shared} work range inside the accepted
domain, not a new result for those illustrative coefficients. Predecision
histories stay unchanged since no traffic reached the terminal.

\subsection{Preparation and its nonzero asymmetric speed sensitivity}
With $z=vt$ and $\Delta_i(v)=\kc(v)-r_i>0$, H1 is
\begin{equation}
 e_{i,z}=-\frac1h\min\{e_i,(\Delta_i+e_i)/2\}.
 \label{eq:h1}
\end{equation}
$\mathcal L$ is donor-limited at every advice: $\Delta_{\mathcal L}\geq
7\kappa/10>e_0$. Its density at every fixed progress is independent of
advice. $\mathcal H$ starts receiving-limited only when
$v>v_c=58u/31$. Its switch progress is
\begin{equation}
 z_s(v)=2h\log\frac{e_0+\Delta_{\mathcal H}(v)}{2\Delta_{\mathcal H}(v)}.
 \label{eq:switch}
\end{equation}
At its largest value, $v=2u$, $\exp(z_s/h)=(23/22)^2<6/5=\exp(\ell/h)$.
Thus all terminal handoff exposures $z\in[\ell,2\ell]$ are donor-branch
exposures, with front density $B_i(v)=e_i(\ell;v)$:
\begin{equation}
 B_{\mathcal L}=\frac\kappa2,\qquad
 B_{\mathcal H}(v)=
 \begin{cases}
 \kappa/2,&v\leq58u/31,\\[2pt]
 \displaystyle\frac56\frac{(e_0+\Delta_{\mathcal H})^2}{4\Delta_{\mathcal H}},
 &v>58u/31.
 \end{cases}\label{eq:front}
\end{equation}
The two expressions match with zero first derivative at $v_c$. At $v_+$,
\begin{equation}
 B_{\mathcal H}=\frac{529\kappa}{1056},\quad
 B'_{\mathcal H}=\frac{115\kappa}{7744u}>0,\qquad
 B_{\mathcal L}=\frac\kappa2,\quad B'_{\mathcal L}=0.
 \label{eq:frontvalues}
\end{equation}
This replaces the historical symmetric-member calculation; its different E
population is not reused as the present hidden variable.

\section{A nonzero common-fastest range in the selected law}
The front free-work test is $W_{f,i}(v)=v[r_i+\delta B_i(v)]$.
H1 comparison makes $B_i$ nondecreasing in advice, so the largest front work
is at $2u$. There,
\begin{equation}
 \frac{W_{f,\mathcal L}}{\kappa u}=\frac85+\delta,\qquad
 \frac{W_{f,\mathcal H}}{\kappa u}=\frac{19}{10}+\frac{529}{528}\delta.
 \label{eq:frontwork}
\end{equation}
Consequently, for
\begin{equation}
 0\leq\delta\leq\frac{132}{2645},\qquad
 J_i(v)=\frac{M(L-\ell/2)}v\quad\hbox{for both states and every advice}.
 \label{eq:negative}
\end{equation}
Joint free admission holds; its workload decreases through the pulse.
Free E's mark is $v\geq U(R)$, so no recovery loss is deposited.
At equality the selected free-continuation guard applies. This includes
positive maneuver burden and can have positive $\beta$ and finite recovery
acceleration. Both preparations uniquely prefer $v_+$. Information can be
missing without being needed for the best action in this range.

\section{Complete TST can have different useful advice changes}
\subsection{An exact auxiliary reference, with every waiting tail counted}
For this subsection only, delete mismatch and retained recovery while
keeping H1, propagation, $c_0$, initial state and TST. Denote this fixed-work
reference by $J_i^0$. A remains free because $\lambda_A< R$; while E queues,
$g_E=(R-\lambda_A)/c_0$. This is an auxiliary reference/rival, not an
amendment of the selected law.

Let $t_f=(L-\ell)/v$, $d=\ell/v$, and $s=t-t_f$ during the pulse.
Since every exposure is past the applicable preparation switch,
\begin{equation}
 \lambda_E=vB_i(v)e^{-vs/h},\quad
 \lambda_A=v[r_i-B_i(v)e^{-vs/h}],\quad 0<s<d.
 \label{eq:arrivalsreference}
\end{equation}
Set
\begin{equation}
 m_i=hB_i(v)(1-e^{-\nu})=hB_i(v)/6,\quad
 Q_i=hB_i(v)\ell-hm_i,\quad D_i=\ell r_i-Rd,\quad
 \mathcal T=M(L-\ell/2).
 \label{eq:referencequantities}
\end{equation}
$m_i$ is forced E mass, distinct from the fixed renewal mass $m$;
$Q_i$ has units mass times length and is an arrival-pulse area quantity.
No normalization or packet size is changed.

A sufficient condition for a strictly positive E queue throughout the
eligible pulse at every advice is
\begin{equation}
 c_0>94/25.
 \label{eq:sustained}
\end{equation}
Indeed $e_i$ at any handoff exposure is at least
$e_0e^{-2\nu}=5\kappa/12$, $r_i\geq4\kappa/5$, $v\geq u$, and
$u[4\kappa/5+(c_0-1)5\kappa/12]>R$. The queue is
\begin{equation}
 q_i^0(s)=\frac{(vr_i-R)s+(c_0-1)hB_i(1-e^{-vs/h})}{c_0}.
 \label{eq:referencequeue}
\end{equation}
It drains after the rear at $R/c_0$. Integrating the pulse and full tail gives
\begin{equation}
 J_i^0(v)=\frac{\mathcal T}{v}+\frac{D_id}{2c_0}
       +(1-1/c_0)\frac{Q_i}{v}
       +\frac{[D_i+(c_0-1)m_i]^2}{2c_0R}.
 \label{eq:referenceobjective}
\end{equation}
Every initial lane's travel and the entire E tail are included; A has no
waiting in this reference. Equation~\eqref{eq:referenceobjective} is used
only on its proved branch.

\subsection{The reference action contrast is proved anew for fixed E}
For $\mathcal L$, $B,m_i,Q_i$ are action-independent. Writing
$W_i=\ell r_i+(c_0-1)m_i$, exact cancellation of the $v^{-2}$ terms gives
\begin{equation}
 J_{\mathcal L}^0(v)=\frac{W_{\mathcal L}^2}{2c_0R}
        +\frac{A_{\mathcal L}}v,\quad
 A_{\mathcal L}=\mathcal T+(1-1/c_0)(Q_{\mathcal L}-\ell m_{\mathcal L})
                      -\frac{\ell^2r_{\mathcal L}}{2c_0}>0.
 \label{eq:lowmonotone}
\end{equation}
To see strictness, $0\leq Q_i\leq\ell m_i$,
$m_i\leq\ell r_i$ and $\mathcal T=(5/2)\ell M>\ell^2r_i$ imply
$A_{\mathcal L}\geq \mathcal T-\ell^2r_{\mathcal L}>0$.
Hence $v_+$ is the unique optimum of the low-A reference.

For $\mathcal H$, expand the same complete formula as
\begin{equation}
 \begin{gathered}
 J_i^0=\frac{c_0m_i^2}{2R}+I_{0,i}+I_{1,i}/c_0,\\
 I_{0,i}=\frac{\mathcal T+Q_i}{v}+\frac{(D_i-m_i)m_i}{R},\quad
 I_{1,i}=\frac{D_id}{2}-\frac{Q_i}{v}+\frac{(D_i-m_i)^2}{2R}.
 \end{gathered}\label{eq:expansion}
\end{equation}
At $v_+$, $A_+=m_{\mathcal H}m'_{\mathcal H}/R>0$, whereas
$I'_{0,\mathcal H}$ and $I'_{1,\mathcal H}$ are finite and independent of
$c_0$. Thus
\begin{equation}
 c_0>\max\left\{94/25,
 \frac{|I'_{0,\mathcal H}(v_+)|+|I'_{1,\mathcal H}(v_+)|}{A_+}\right\}
 \quad\Longrightarrow\quad (J_{\mathcal H}^0)'(v_+)>0.
 \label{eq:positivecost}
\end{equation}
Near $v_-$ the high-A preparation is donor-only, and the same positive
$A_i/v$ coefficient gives a strictly negative derivative. The continuous
reference therefore has interior high-A minimizers and disjoint low/high
minimizing sets. A sufficiently small common feasible slowdown from $v_+$
strictly worsens $\mathcal L$ and improves $\mathcal H$. No physical
plausibility or moderate-work result follows from this existence bound.

\subsection{Transfer finite action utility into the selected positive-mechanism law}
Restore the actual fixed W-F mismatch, E-only source and recent-mass O-R
renewal. With $u_R=U(R)$, its zero-mark isolated loss area and account satisfy
\begin{equation}
 W_0=\frac{K}{w+u_R}\frac{u_R^2}{2a_r},\qquad
 0\leq p(t)\leq\varepsilon_p:=\frac{RW_0}{mc_0}.
 \label{eq:lossbound}
\end{equation}
The latter bounds the nonnegative retained-window time integral by the
all-history integral with $g_E\leq R/c_0$. It does not replace renewal.
Put $a_*=v_+\max_i r_i<R$ and choose
\begin{equation}
 0<\beta\leq1,\quad c_0-\beta(1+v_+/w)>0,\quad
 \varepsilon_p<\min\{(R-a_*)/2,R/2\}.
 \label{eq:smallparameters}
\end{equation}
Then $\Bud\geq R-\varepsilon_p>a_*$, so A is free by its actual admission
test. While E queues, its service is bracketed by
\begin{equation}
 \mu_-:=\frac{R-\varepsilon_p-\lambda_A}{c_0+\beta}
 \leq g_E\leq\mu_0:=\frac{R-\lambda_A}{c_0},\quad
 0\leq\mu_0-\mu_-\leq\varepsilon_\mu
 :=\frac{\varepsilon_p}{c_0}+\frac{R\beta}{c_0^2}.
 \label{eq:servicebrackets}
\end{equation}
The same arrivals drive reflected comparison queues $q_0,q_-$ with these
two positive services. The causal free/queued guards and a first-crossing
comparison give $q_0\leq q_E\leq q_-$; the actual queue itself need not be a
fixed-service reflection. Reflection implies $q_--q_0\leq t\varepsilon_\mu$.
Choose $H>L/v_-+2(c_0+1)M/R$ after fixing $c_0$ and geometry. Every compared
cohort clears by $H$. Consequently, for both states and all advice,
\begin{equation}
 0\leq J_i(v)-J_i^0(v)\leq\varepsilon_J:=H^2\varepsilon_\mu/2.
 \label{eq:valuebound}
\end{equation}
The initial split is held fixed in this estimate. It makes no derivative
closeness assertion. With $c_0$ fixed first, positive $\beta$ can be made
small and positive finite $a_r$ sufficiently large to make $\varepsilon_J$
as small as needed while retaining every selected mechanism.

Let $d_i^0(v)=J_i^0(v)-\min_zJ_i^0(z)$. Compactness, continuity and
disjoint reference minimizing sets give
\begin{equation}
 \Gamma_0=\min_v\max_i d_i^0(v)>0.
 \label{eq:referenceregret}
\end{equation}
For actual $V_i=\inf_vJ_i(v)$, uniform comparison yields conservatively
\begin{equation}
 \inf_v\max_i\{J_i(v)-V_i\}\geq\Gamma_0-2\varepsilon_J>0
 \qquad(\varepsilon_J<\Gamma_0/4).
 \label{eq:actualregret}
\end{equation}
No global continuity or attainment of actual minima is assumed. Making
$2\varepsilon_J$ smaller than each of the two strict reference slowdown
gaps also preserves their opposite signs in the selected law. Thus identical
Level 2 histories can require incompatible near-optimal common advice.
Level 3 and its compact peer can distinguish the states.

\textbf{Limits.} This is a new asymmetric action certificate, not the old
hidden-E result renamed. It may need large work, weak mismatch and fast
recovery. A stays free in this range; the archived A-delay/active-retirement
witness is separate. No joint strong-delay/action/physical witness, default
calibration, opposite infinitesimal gradient signs in the actual law,
AD advantage or real-freeway usefulness is proved.

\section{The complete sensitivity reference does not require initial symmetry}
\subsection{Preparation, clocks and all three lane arrivals}
Define $E_v(z;v)=\partial_ve(z;v)$ at fixed physical progress, with
$E_v(0;v)=0$. On each smooth H1 branch,
\begin{equation}
 (E_v)_z=\begin{cases}-E_v/h,&\text{donor},\\
 -(E_v+\kc'(v))/(\gamma h),&\text{receiving},\end{cases}
 \quad\kc'(v)=-K/(w+v)^2.
 \label{eq:prepsensitivity}
\end{equation}
The vector fields match at their moving switch, so the state sensitivity
has no jump from unequal vector fields there. Every true $r,b_0,e_0$ is
fixed within the advice derivative; one does not regenerate initial traffic
or recondition its history after changing advice.

For an initial position $x$, eligibility is $(L-x)/v$. Exactly,
\begin{equation}
 P_E(t;v)=\int_0^\ell e(L_p-x;v)\mathbf1_{(L-x)/v\leq t}\dd x.
 \label{eq:PE}
\end{equation}
Inside the eligible pulse set $x_*=L-vt$, $z_*=vt-b$. Then
\begin{align}
 P_{E,v}&=\int_{x_*}^\ell E_v(L_p-x;v)\dd x+t e(z_*;v),\nonumber\\
 \lambda_E&=ve(z_*;v),\qquad
 \lambda_{E,v}=e+vE_v+vt e_z,\nonumber\\
 P_{A,v}&=rt-P_{E,v},\quad P_{B,v}=b_0t,\quad
 \lambda_{A,v}=r-\lambda_{E,v},\quad\lambda_{B,v}=b_0,\nonumber\\
 t_f'&=-(L-\ell)/v^2,\qquad t_r'=-L/v^2.
 \label{eq:arrivalsensitivity}
\end{align}
Before the front all cumulative sensitivities vanish; after the rear the
E derivative is its full progress integral, A is its negative and B is zero.
Their final total sensitivity is zero by fixed total mass. Use one-sided
front/rear values; fixed clocks would omit a genuine causal channel.

\subsection{Implicit point service and its traffic marks}
For queued E retain $\He(g;a,u)=g[c_0+\chi a\pos{u-U(g)}]=x$.
Here $a$ is A flow. On positive mismatch,
\begin{equation}
 \begin{gathered}
 \He_g=c_0+\chi a\{u-U(g)-gU'(g)\}>0,\quad
 \He_a=\chi g\{u-U(g)\},\quad\He_u=\chi ag,\\
 g_v=\frac{x_v-\He_a a_v-\He_u u_v}{\He_g}.
 \end{gathered}\label{eq:rootsensitivity}
\end{equation}
On inactive mismatch, $\He_g=c_0$, $\He_a=\He_u=0$.
Free A/queued E uses $(a,u,x)=(\lambda_A,v,\Bud-\lambda_A)$;
both queued use $(\Bud/2,U(\Bud/2),\Bud/2)$.
Free marks have derivative one; queued marks have derivative $U'(g)g_v$.
The other service branches are differentiated from their retained admission
equations. At fixed current traffic/budget the fully queued root has no
direct advice term. $\chi$ is not recomputed with advice or a search box.

\subsection{History, retirement and queue events}
Write $Y=Y_E$, $g=g_E$ and $Z=Y_v$. On a regular history chart,
\begin{equation}
 p(t)=\frac1m\int_{s_m(t)}^tD(t-s;u_E(s))g_E(s)\dd s,
 \quad Y(s_m)=Y(t)-m\quad\hbox{after fill},
 \label{eq:history}
\end{equation}
with $s_m=0$ before fill. After fill suppose the cutoff lies at positive
old E rate, away from a mark event. Then
$g(s_m)s_{m,v}=Z(t)-Z(s_m)$. For source/mark events $\tau_k(v)$ inside
the retained window put
$\Phi_k^\pm=D(t-\tau_k;u_E(\tau_k^\pm))g_E(\tau_k^\pm)$ and
$D_o=D(t-s_m;u_E(s_m))$. Piecewise Leibniz differentiation gives
\begin{align}
 p_v=\frac1m\bigg[&\int_{s_m}^t(D_u u_{E,v}g_E+Dg_{E,v})\dd s
 +\sum_k(\Phi_k^- -\Phi_k^+)\tau_k'\nonumber\\
 &-\mathbf1_{Y(t)>m}D_o\{Z(t)-Z(s_m)\}\bigg],\qquad \Bud_v=-p_v.
 \label{eq:historysensitivity}
\end{align}
The last term is old-weight retirement; the old-rate denominator cancels.
It must not be omitted or replaced by a strictest-old-cap rule. For active
loss, with $S=\sqrt{(w+u)^2+2wa_r\theta}$,
\begin{equation}
 D_u=-Kw(w+u)/S^3,\qquad D_\theta=-Kw^2a_r/S^3.
 \label{eq:kernel}
\end{equation}
They are zero on the inactive side. At fixed integration time $s$, physical
age $t-s$ is fixed; moving source-event times enter through the jump terms.
The zero-loss support edge has no jump contribution because $D$ is continuous
and zero there. Fill, plateau and cutoff/mark encounters require their own
one-sided/event treatment, not a universal smooth derivative assertion.

For an isolated transverse guard let $G(t,v)$ denote its left-side value
evaluated along the full history-aware solution. Then
\begin{equation}
 \tau_k'=-G_v^-/G_t^-,\qquad
 X_v^+=X_v^-+(F^- -F^+)\tau_k'
 \label{eq:events}
\end{equation}
for continuous physical-state coordinates with identity reset. Guard
derivatives include the history sensitivity above. This notation does not
claim a finite-dimensional Markov reduction of the account. A birth uses
the free-feasibility margin $\Bud-\lambda_A$, not the identically zero
free A queue; clearance uses the positive queue approaching zero. Recorded
source-event times remain in the history. The regular timing principle is
the one reviewed in Kong et al.; the renewal derivative is a separate
model-specific deduction.

\subsection{Complete TST and the derivative's domain}
Use a fixed horizon above the AS1 clearance bound, independent of advice:
\[
 H_*>L/v_-+c_{\max}M/B_{\min},\qquad B_{\min}>0.
\]
With $Z_j=Y_{j,v}$,
\begin{align}
 J&=\int_0^{H_*}(M-Y_A-Y_E-Y_B)\dd t
   =\frac{M(L-\ell/2)}v+\int_0^{H_*}(q_A+q_E)\dd t,\nonumber\\
 J'&=-\int_0^{H_*}(Z_A+Z_E+Z_B)\dd t
   =-\frac{M(L-\ell/2)}{v^2}+\int_0^{H_*}(q_{A,v}+q_{E,v})\dd t.
 \label{eq:tstgradient}
\end{align}
Queue-area endpoint values vanish at zero-queue hits; their timing effects
on future service/history do not. All traffic, including both A/E waiting
tails, is counted. Recovery after the last completed unit is not extra own
travel beyond the selected choke completion.

On a neighborhood with a fixed finite event list, smooth branches, transverse
guards and regular cutoffs, differentiating the locally uniform causal
Volterra contraction gives the within-chart linear sensitivity. Implicit
guard differentiation and identity-state matching join the charts; bounded
regular-chart terms permit fixed-horizon differentiation. This proves the
local reference construction. It uses no initial A/B equality.

Forward-admitted grazing, coincident events, activation thresholds, root/
profile kinks and renewal encounters remain in the task. Check continuity
and feasible side behavior before assigning a derivative. If no finite
ordinary or one-sided derivative exists, report that and the actual finite
cost change. No smoothing, frozen event, deleted boundary case or new
generalized-gradient controller is selected. Uniform value comparison does
not imply gradient closeness or global minimum attainment.

\section{Evidence and source limits at the original Step 3 gate}
AS1's complete-forward result, AS2's conditional information contrast and
compact equality, and this asymmetric action/sensitivity review now fit one
standalone verification contract. Later evidence must distinguish continuous
reference correctness, numerical forward convergence, numerical differences,
AD, feasible-step usefulness and optimizer performance. A strong bounded
scalar search and the compact peer remain required rivals. No such numerical
layer is marked passed here.

Information revealed through complicated features may be harder for a
learner to extract; that is a plausible separate conditioning/learning
hypothesis. AS2's exactly identical histories establish structural ambiguity,
not that learning hypothesis. When information is recoverable from the same
inputs, model-aware inference and feature transforms must be allowed equally;
learning difficulty needs matched tasks, inputs and resources before a claim.
No training experiment is opened by this clarification.

This note adapts the own fixed-work derivation and service/queue comparison
in the historical
\href{vsl_step3_standalone_information_and_gradient.pdf}{standalone Step 3 mathematics},
with the new fixed-E split calculation and explicit sensitivity audit above.
It does not transplant the historical symmetric action certificate.
The archived \href{https://arxiv.org/pdf/2306.06862}{Kong et al. hybrid-sensitivity paper}
supports regular event-time corrections only; the
\href{https://arxiv.org/pdf/1603.04785}{Delle Monache--Piccoli--Rossi VSL paper}
is prior art under a different objective. No paper is newly viewed or new
traffic model adopted in AS3. Applied MP6 remains unexecuted, SUMO optional,
Step 4/implementation unopened and ramp/urban lines closed.
\end{document}
